\documentclass[11pt]{article}

\usepackage[margin=1in]{geometry}
\usepackage{amsmath,amssymb}
\usepackage{booktabs}
\usepackage{tabularx}
\usepackage{array}
\usepackage{enumitem}
\usepackage[T1]{fontenc}
\usepackage{hyperref}

\newcommand{\code}[1]{\texttt{#1}}
\newcommand{\Uh}{\hat U}
\newcommand{\Dc}{D_{\mathrm{crit}}}
\newcommand{\Dsn}{D_{\mathrm{SN}}}
\newcolumntype{L}{>{\raggedright\arraybackslash}X}
\emergencystretch=3em

\title{Cluster crystals near the instability threshold:\\
what the 2D and 3D \code{gema} runs show and how to interpret it}
\author{FHDeX}
\date{}

\begin{document}
\maketitle

\noindent This note collects the analysis of the repulsive GEM-$\alpha$
(\code{gema}) runs in \code{exec/dean\_kow/interaction}, in 2D and 3D, near the
linear instability of the uniform state. It covers what was observed, the
theory that explains it, and how to read the diagnostics. The scripts are
described in \code{README\_crystal\_analysis.md}; the commands that reproduce
each result are in Section~\ref{sec:repro}.

\paragraph{Summary.}
\begin{itemize}[nosep]
  \item The linear threshold $\Dc$ printed by the code is exact for the
        discretized interaction. Noise does not move it.
  \item The transition to the cluster crystal is first-order (transcritical) in
        both 2D and 3D. The cause is the resonant \emph{triads} of the quadratic
        nonlinearity.
  \item \textbf{2D.} Hexagonal crystals exist up to $\Dsn\approx1.18\,\Dc$.
        Near $\Dc$ they form much faster than linear theory predicts. Above
        $\Dc$ they can nucleate from the stable uniform state, and a
        crystal restarted above $\Dc$ survives.
  \item \textbf{3D.} bcc is the preferred cluster crystal, and survives up to
        $\Dsn\approx1.48\,\Dc$.
  \item \textbf{3D quenches.} A deep quench ($D=0.83\,\Dc$) gives an
        amorphous cluster packing. A shallow quench ($D\approx\Dc$) gives a
        box-aligned bcc crystal with twins.
\end{itemize}

\section{Model, parameters and linear theory}

The density $\phi$, normalized to $\int\phi=1$ on the unit periodic box
($\mu_0=1$), obeys the Dean--Kawasaki equation
\begin{equation}
  \partial_t\phi = \nabla\cdot\bigl(D\nabla\phi + \phi\nabla(U\star\phi)\bigr)
     + N^{-1/2}\nabla\cdot\bigl(\sqrt{2D\phi}\,\boldsymbol\xi\bigr),
  \qquad U = \varepsilon\,e^{-(r/R)^\alpha} .
\end{equation}
Linearizing about $\phi=\mu_0$ gives the growth rate of mode $\mathbf k$,
\begin{equation}
  \sigma(k) = -k^2\bigl(D + \mu_0\Uh(k)\bigr), \qquad
  \Lambda = -\mu_0\Uh_{\min}/D, \qquad \Dc = \mu_0|\Uh_{\min}| ,
\end{equation}
and the uniform state is unstable when $\Lambda>1$. In the code, $\Uh$ is the
unnormalized DFT of the sampled potential. The factor \code{cellvol} that
multiplies it in the printout is the integration weight, so
$\Uh\cdot\code{cellvol}$ is the continuum transform (Riemann sum) and also
the exact symbol of the convolution the code applies.
\begin{center}
\begin{tabular}{@{}lllllll@{}}
\toprule
& $\varepsilon$ & $R$, $\alpha$ & grid & $k_c$ & critical box mode & $\Dc$ \\
\midrule
2D (\code{inputs\_fv}) & 666 & 0.1, 3 & $128^2$ & 49.7 (cont.) & $(8,0)$, $|k|=16\pi$ & 0.5472 \\
3D (\code{inputs\_fv\_3d}) & 4509 & 0.1, 3 & $128^3$ & 53.7 & $(3,0,8)$, $|k|=53.68$ & 0.3012 \\
\bottomrule
\end{tabular}
\end{center}
The 3D $\varepsilon$ is chosen so that $\Uh(0)=18.89$ is the same as in 2D.

\paragraph{Diagnostics used.}
\begin{itemize}
  \item \textbf{Structure factor.} $S(\mathbf k)=|\sum n\,e^{-i\mathbf k\cdot
        \mathbf x}|^2/N$ with $n=N\phi/\sum\phi$, so a Poisson field has $S=1$.
        It is recorded on the critical shell against time
        (\code{sk\_growth.py}). The time-averaged full $S(\mathbf k)$ gives the
        spot statistics (\code{sf\_spots.py}): max/mean on the peak shell, the
        share of the brightest modes, the participation ratio
        $\mathrm{PR}=(\sum S)^2/(n\sum S^2)$, and the angles between the
        brightest spots.
  \item \textbf{Bond order of the cluster centres}, from the cluster tools
        (see \code{order\_diagnostics.tex}): $|\psi_6|$ in 2D, Steinhardt
        $q_4,q_6$ (local) and $Q_4,Q_6$ (global) in 3D. Reference values come
        from \code{order\_reference.py}.
\end{itemize}

\section{Weakly nonlinear theory: triads and hexagons above threshold}\label{sec:wnl}

\subsection{Triads}
Write $\phi=\mu_0(1+u)$. The only nonlinearity is
$\mu_0\nabla\cdot(u\nabla(U\star u))$, with Fourier form
\begin{equation}
  N_{\mathbf k} = \sum_{\mathbf p+\mathbf q=\mathbf k} B(\mathbf p,\mathbf q)\,u_{\mathbf p}u_{\mathbf q},
  \qquad
  B(\mathbf p,\mathbf q) = -\frac{\mu_0}{2}\Bigl[(\mathbf k\cdot\mathbf q)\Uh(q)
     + (\mathbf k\cdot\mathbf p)\Uh(p)\Bigr].
\end{equation}
The product of modes $\mathbf p$ and $\mathbf q$ drives the mode
$\mathbf p+\mathbf q$. This is three-wave mixing, as in nonlinear optics.
\begin{itemize}
  \item \textbf{Non-resonant pairs.} Usually $\mathbf p+\mathbf q$ is off the
        critical shell. It is then damped and only feeds back at cubic order.
  \item \textbf{Triads.} A \emph{triad} is three critical wavevectors with
        $\mathbf k_1+\mathbf k_2+\mathbf k_3=0$. Since all have length $k_c$,
        they form an equilateral triangle, each pair $120^\circ$ apart. For a
        triad the interference pattern of two waves drives the third directly,
        giving a \emph{quadratic} term in the amplitude equation.
  \item \textbf{Phase.} Translations shift each wave's phase but not the sum
        $\Phi=\theta_1+\theta_2+\theta_3$. The sign of the coupling selects
        $\Phi=0$, so the three cosines peak together at lattice points:
        clusters, not a honeycomb of holes.
  \item \textbf{Landau theory.} A triad allows a cubic invariant
        $-\gamma A_1A_2A_3\cos\Phi$ in the Landau free energy, and a cubic
        invariant forces a first-order transition.
\end{itemize}
\begin{center}
\begin{tabular}{@{}lll@{}}
\toprule
pattern & critical vectors & triads \\
\midrule
stripes, lamellae & $\pm\mathbf k$ & none ($\mathbf k-\mathbf k=0$ is the conserved mass) \\
2D squares, 3D simple cubic & $90^\circ$ apart & none \\
2D hexagons, 3D hexagonal rods & 6 at $60^\circ$ & 2; each vector is in 1 \\
3D fcc & 8 $\langle111\rangle$ & none on the first shell \\
3D bcc & 12 $\langle110\rangle$ & 8; each vector is in 2 \\
\bottomrule
\end{tabular}
\end{center}

\subsection{Amplitude equations}
Put amplitudes $A_j$ on the critical modes and slave every other mode the
quadratic term generates, $w_{\mathbf K}=-N_{\mathbf K}/\sigma(K)$. For 2D
hexagons, with $\Uh_c$, $\Uh_{\sqrt3}$ and $\Uh_2$ the values of $\Uh$ at
$k_c$, $\sqrt3k_c$ and $2k_c$:
\begin{gather}
  \dot A_1 = \sigma A_1 + \gamma A_2^*A_3^* - \bigl[g|A_1|^2 + h(|A_2|^2+|A_3|^2)\bigr]A_1, \\
  \gamma = k_c^2\mu_0|\Uh_c|, \quad
  g = \frac{k_c^2\mu_0\Uh_c(\mu_0\Uh_c/2-\mu_0\Uh_2)}{D+\mu_0\Uh_2}, \quad
  h = \frac{k_c^2\mu_0\Uh_c(\mu_0\Uh_c-3\mu_0\Uh_{\sqrt3})}{2(D+\mu_0\Uh_{\sqrt3})},
\end{gather}
with $\sigma=k_c^2(\Dc-D)$. For any lattice with equal amplitudes the same
reduction (\code{amplitude\_eqs.py}) gives
\begin{equation}
  \dot A = \sigma A + \gamma_{\mathrm{eff}}A^2 - GA^3, \qquad
  \Dsn = \Dc + \frac{\gamma_{\mathrm{eff}}^2}{4Gk_c^2}, \qquad
  A(\Dc) = \gamma_{\mathrm{eff}}/G .
\end{equation}
The lattice exists for $D<\Dsn$, a saddle-node. Between $\Dc$ and $\Dsn$ both
it and the uniform state are linearly stable. The unstable branch between them,
\begin{equation}
  A_u = \frac{\gamma_{\mathrm{eff}} - \sqrt{\gamma_{\mathrm{eff}}^2+4\sigma G}}{2G}
      \approx \frac{|\sigma|}{\gamma_{\mathrm{eff}}} ,
\end{equation}
is the barrier noise must cross. When $\Uh$ is negligible at $\sqrt3k_c$ and
$2k_c$, $g=h=\gamma/2$ and $\Dsn=\tfrac76\Dc$ for 2D hexagons, whatever the
kernel. For $0<\sigma\ll\gamma$, the quadratic term dominates once
$A>A_c=\sigma/\gamma$. After that, $\dot A\approx\gamma A^2$ runs away in a
time of about $1/(\gamma A_0)$, independent of $\sigma$.

\begin{center}
\begin{tabular}{@{}lllllll@{}}
\toprule
& lattice & $\gamma_{\mathrm{eff}}$ & $G$ & $A(\Dc)$ & $\Dsn/\Dc$ & note \\
\midrule
2D & hexagons & 1383 & 2086 & 0.66 & 1.166 & $g=687$, $h=700$ \\
2D & stripes, squares & 0 & 687, 446 & --- & 1 & supercritical \\
3D & bcc & 1737 & 2035 & 0.85 & 1.427 & 8 triads \\
3D & hexagonal rods & 868 & 1310 & 0.66 & 1.166 & \\
3D & fcc & 0 & $-9142$ & --- & --- & one-shell theory fails \\
3D & sc, lamellae & 0 & 280, 433 & --- & 1 & supercritical \\
\bottomrule
\end{tabular}
\end{center}
The 2D values are at the box mode $(8,0)$ with the code's discrete
$\Dc=0.5472$; the 3D values are continuum. For fcc, the $\langle200\rangle$ mode
generated by $\langle111\rangle+\langle11\bar1\rangle$ lies at $1.155\,k_c$,
where $\mu_0\Uh=-0.189$, 63\% of the minimum. That mode is nearly unstable, so
it enters through a small denominator and drives $G$ negative. fcc needs both
shells active, and only the full equation settles it.

\section{The full deterministic equation}\label{sec:cont}
\code{lattice\_continuation.py} solves the noise-free equation
pseudo-spectrally on one periodic cell of a lattice: 2 clusters for hexagons
and bcc, 4 for fcc. The potential is periodized over the cell, so the cell's
Fourier modes see the continuum $\Uh$. Each $D$ is run in turn, starting from
the previous state, which follows the stable branch to its saddle-node. The
free energy
$\mathcal F=\int[D\phi\ln\phi+\tfrac12\phi\,U\star\phi]$ is compared with the
uniform value $\tfrac12\mu_0^2\Uh(0)$.

\paragraph{Checks of the amplitude equations (2D).}
\begin{itemize}
  \item \textbf{Stripes.} At small $\sigma>0$, $A=\sqrt{\sigma/g}$ with a
        measured $g=688$--$697$, against the predicted 687.
  \item \textbf{Hexagons at $\Dc$.} $A=0.71$, against $\gamma/G=0.66$.
\end{itemize}

\begin{center}
\begin{tabular}{@{}lllll@{\qquad}lllll@{}}
\toprule
\multicolumn{5}{c}{2D hexagons ($|\mathbf k|=16\pi$)} & \multicolumn{5}{c}{3D, $\Delta\mathcal F$ per volume} \\
$D$ & $D/\Dc$ & $\phi_{\max}$ & $A$ & $\Delta\mathcal F$ & $D$ & $D/\Dc$ & bcc $\phi_{\max}$ & bcc & fcc ($\times1.05$) \\
\midrule
0.50 & 0.914 & 16.0 & 0.764 & $-0.166$ & 0.25 & 0.83 & 187 & $-0.508$ & $-0.489$ \\
0.55 & 1.005 & 13.0 & 0.709 & $-0.093$ & 0.30 & 1.00 & 128 & $-0.338$ & $-0.317$ \\
0.60 & 1.096 & 9.98 & 0.627 & $-0.033$ & 0.36 & 1.20 & 82 & $-0.163$ & $-0.138$ \\
0.62 & 1.133 & 8.62 & 0.579 & $-0.015$ & 0.39 & 1.30 & 65 & $-0.088$ & $-0.062$ \\
0.63 & 1.151 & 7.85 & 0.547 & $-0.006$ & 0.42 & 1.39 & 48 & $-0.023$ & $+0.004$ \\
0.64 & 1.170 & 6.94 & 0.505 & $+0.001$ & 0.43 & 1.43 & 42 & $-0.004$ & \\
0.6475 & 1.183 & 6.01 & 0.456 & $+0.005$ & 0.44 & 1.46 & 36 & $+0.013$ & \\
0.65 & 1.188 & \multicolumn{3}{l}{decays} & 0.45 & 1.49 & \multicolumn{3}{l}{bcc decays} \\
\bottomrule
\end{tabular}
\end{center}
\begin{center}
\begin{tabular}{@{}llll@{}}
\toprule
& $\Dsn$ (full equation) & $\Dsn$ (amplitude eqs.) & $D_F$ ($\Delta\mathcal F=0$) \\
\midrule
2D hexagons & $0.648\pm0.002 = 1.18\,\Dc$ & 1.17 & $0.639 = 1.17\,\Dc$ \\
3D bcc & $0.445\pm0.005 = 1.48\,\Dc$ & 1.43 & $0.432 = 1.44\,\Dc$ \\
3D fcc ($\times1.05$) & $>0.42$ & --- & $0.418 = 1.39\,\Dc$ \\
\bottomrule
\end{tabular}
\end{center}
\begin{itemize}
  \item \textbf{Free energy.} Below $D_F$ the lattice is the global
        free-energy minimum; between $D_F$ and $\Dsn$ it is metastable.
  \item \textbf{fcc spacing.} fcc prefers a spacing about 5\% larger than
        $|\mathbf G_{111}|=k_c$. At $D=0.30$, scale factors 0.95, 1.00, 1.05
        and 1.10 give $\Delta\mathcal F=-0.12$, $-0.28$, $-0.32$ and $-0.20$.
        bcc is best at scale 1.00.
  \item \textbf{bcc against fcc.} bcc is lower than fcc at every $D$ tested, by
        about 0.02. This is the Alexander--McTague preference for bcc near a
        weakly first-order freezing.
  \item \textbf{Low-$D$ accuracy.} At $D=0.25$ the 3D clusters are about two
        grid cells wide, so the low-$D$ free energies are the least accurate.
\end{itemize}

\section{2D runs}

\subsection{Below threshold: noise-driven, then triad-driven onset}
Five runs ($128^2$, uniform start) bracket $\Dc$ from below. A sixth,
\code{\_new\_2d\_10\_sdt}, repeats $D=0.46$ with a smaller time step and gives
the same result. ``Onset'' is when the shell mean of $S$ reaches $10^4$; the
linear prediction integrates $\dot S=2\sigma S+2k^2D$ from $S=0$.
\begin{center}
\begin{tabular}{@{}lllllllll@{}}
\toprule
run & $D$ & $\sigma$ & $N$/cell & $A_c=\sigma/\gamma$ & rate$/\sigma$ & onset & linear & clusters \\
\midrule
\code{farther} & 0.4165 & 330 & 20 & 0.24 & 0.9 & 0.014 & 0.012 & 56 \\
\code{\_new\_2d\_10} & 0.46 & 220 & 10 & 0.16 & 1.0 & 0.021 & 0.017 & 56 \\
\code{closer} & 0.503 & 112 & 20 & 0.08 & 1.3 & 0.025 & 0.030 & 53 \\
\code{vclose} & 0.54 & 18.2 & 20 & 0.013 & 2.1 & 0.066 & 0.135 & 56 \\
\code{vvclose} & 0.545 & 5.6 & 20 & 0.004 & 7.6 & 0.061 & 0.334 & 56 \\
\bottomrule
\end{tabular}
\end{center}
\textbf{What is seen.} Far from $\Dc$, growth follows linear theory. Close to
$\Dc$ it is several times faster, and the two closest runs form clusters at
nearly the same time even though $\sigma$ differs by $3\times$.

\textbf{Interpretation.} Noise first builds $S$ to about 20 in $t\approx0.01$,
a per-mode amplitude $A=\sqrt{S/N}\approx0.01$, nearly independent of $\sigma$.
\begin{itemize}
  \item \textbf{Far from $\Dc$} ($A_c\gg0.01$), the modes grow linearly almost
        to saturation.
  \item \textbf{Near $\Dc$} ($A_c\lesssim0.01$), the triad term takes over at
        once. Onset is then set by $1/(\gamma A_0)\approx0.07$, not by
        $1/\sigma$. The formula $t_{\mathrm{form}}\approx\ln N/2\sigma$ fails there.
\end{itemize}
In every run the final state is a hexagonal crystal of 53--56 clusters, set by
the $(8,0)$ mode.

\subsection{Above threshold: nucleation and hysteresis}
For $D>\Dc$ the fluid is linearly stable, and the critical modes fluctuate
about $S_{\mathrm{eq}}=1/(1-\Lambda)$. That is a per-mode noise amplitude
$A_{\mathrm{noise}}\approx[N(1-\Lambda)]^{-1/2}$, to be compared with the
barrier $A_u$.
\begin{center}
\begin{tabularx}{\textwidth}{@{}lllllllL@{}}
\toprule
run & $D$ & $\Lambda$ & $S_{\mathrm{eq}}$ & $A_{\mathrm{noise}}$ & $A_u$ & $A_u/A_{\mathrm{noise}}$ & observed \\
\midrule
\code{labove} & 0.56 & 0.977 & 44 & 0.016 & 0.024 & 1.5 & $S\approx30$ until $t\approx0.17$, then nucleation; full crystal by $t=0.21$ \\
\code{labove/\_D\_.57} & 0.57 & 0.960 & 25 & 0.012 & 0.045 & 3.6 & no clusters to $t=0.42$; shell $S\approx22$ \\
\code{above} & 0.60 & 0.912 & 11 & 0.008 & 0.12 & 14 & no clusters to $t=0.92$ \\
\code{*/\_det} & 0.56, 0.60 & & & & & & no noise: decays to uniform \\
restart & 0.64 & 0.855 & & & & & hexagonal state restarted at $D=0.64$ keeps its lattice \\
\bottomrule
\end{tabularx}
\end{center}
The \code{above} and \code{labove} runs have $N=163840$ (10 particles per
cell). The \code{\_D\_.57} sub-run has no inputs file and is assumed to match;
its shell $S$ agrees with $S_{\mathrm{eq}}$ for that $N$. The restart at
$D=0.64$ is the hysteresis test reported for these runs; its run directory is
not part of the analysed set.
\begin{itemize}
  \item \textbf{Linear stability holds.} Without noise the uniform state
        recovers. With noise, the shell $S$ sits at about $S_{\mathrm{eq}}$
        until clusters nucleate.
  \item \textbf{Nucleation.} At $D=0.56$ the barrier is only $1.5\times$ the
        noise. After a waiting time of about 0.17 a nucleus forms and grows to
        a single hexagonal crystal (global $|\psi_6|=0.70$, $\mathrm{PR}=6/48$:
        six Bragg spots) in about 0.04. At $3.6\times$ and $14\times$ the
        barrier, no nucleation was seen on the times run.
  \item \textbf{Hysteresis.} A crystal restarted at $D=0.64<\Dsn$ survives,
        as the continuation predicts. A restart above $\Dsn\approx0.65$ should
        melt.
\end{itemize}
This is the behaviour of a first-order transition with spinodal $\Dc$:
nucleation and hysteresis in the window $\Dc<D<\Dsn$.

\section{3D runs}

\subsection{Predictions}
\begin{itemize}
  \item \textbf{Lattice.} bcc has the most triads, the largest window
        ($\Dsn\approx1.48\,\Dc$) and the lowest free energy, so it is the
        mean-field cluster crystal.
  \item \textbf{Box commensurability.} The bcc with reciprocal vectors
        $2\pi(6,6,0)$ and permutations ($|\mathbf k|=53.3$, essentially on the
        critical ring) fits the unit box exactly, with conventional cell
        $1/6$ and 432 clusters.
  \item \textbf{Onset.} With $N=2\cdot10^7$ the noise is weak ($A\approx10^{-3}$
        at $t\approx0.01$). Linear growth therefore carries the amplitude to
        $A_c=\sigma/\gamma$ before the triads take over. The estimated onset
        is about $0.04$ at $D=0.25$, $0.13$ at $0.29$, $0.18$ at $0.295$ and
        $0.25$ at $0.298$.
\end{itemize}

\subsection{Deep and shallow quench}
\begin{center}
\begin{tabular}{@{}lll@{}}
\toprule
& deep quench & shallow quench \\
& $D=0.25$, \code{\_new\_3d\_noise} & $D=0.30$, \code{near\_3d} \\
\midrule
$D/\Dc$, $\sigma$ & 0.83, 149 & 0.996, $\approx3.5$ \\
time analysed & $t=0.040$ & $t=0.171$ \\
clusters & 482 & 457 (89\% of the mass) \\
$q_4$, $q_6$ (12 nn) & 0.127, 0.422 & 0.121, 0.454 \\
$Q_4$, $Q_6$ (12 nn) & 0.013, 0.066 & 0.020, 0.224 \\
$q_4$, $q_6$, $Q_4$, $Q_6$ (cut 0.2) & --- & 0.067, 0.400, 0.005, 0.200 \\
peak shell max/mean & 10.7 & 27.5 \\
brightest-12 share & 0.12 & 0.29 \\
PR & 0.36 & 0.090 \\
full $S(\mathbf k)$ & uniform spherical shell & discrete spots \\
\bottomrule
\end{tabular}
\end{center}
Reference values (\code{order\_reference.py}, 432 points, neighbours within
0.2, columns $q_4$, $q_6$, $Q_4$, $Q_6$):
\begin{itemize}[nosep]
  \item perfect bcc: 0.036, 0.511, 0.036, 0.511;
  \item bcc, 5\% displacement: 0.069, 0.470, 0.038, 0.463;
  \item bcc, 10\% displacement: 0.121, 0.369, 0.024, 0.322;
  \item two random bcc grains: 0.17, 0.32, 0.020, 0.145;
  \item random points (12 nn): $q_6\approx0.27$, $Q_6\approx0.02$.
\end{itemize}

\paragraph{Spots of the shallow quench.} The brightest first-shell modes
($|\mathbf m|=8.49$) are the complete set $(6,6,0)$, $(6,\bar6,0)$, $(6,0,6)$,
$(6,0,\bar6)$, $(0,6,6)$ and $(0,6,\bar6)$, each pair $60^\circ$ or $90^\circ$
apart. Higher shells line up too: $\langle211\rangle$ at
$|\mathbf m|=14.7$ and $\langle220\rangle$ at 17. This is the predicted
box-aligned bcc. A second strong set,
$(2,2,8)$, $(2,8,2)$, $(8,2,2)$, is the image of $(6,6,0)$, $(0,6,6)$,
$(6,0,6)$ under a $180^\circ$ rotation about $[111]$. That rotation leaves
the three vectors perpendicular to $[111]$ unchanged, which is the $\Sigma3$
twin of bcc. Weaker spots at $|\mathbf m|=8.54$, such as $(6,\bar6,\bar1)$,
lie about $7^\circ$ off the main set: mosaic sub-grains. The per-cluster $q_6$
(cut 0.2) has a single broad peak at 0.40, ending at about 0.5, the perfect-bcc
value. So the sample is crystalline throughout, not crystal coexisting with
amorphous regions.

\paragraph{Interpretation.}
\begin{itemize}
  \item \textbf{Deep quench: amorphous.} Local $q_6$ is liquid-like, global
        $Q_6$ is only $3\times$ the random floor, and $S(\mathbf k)$ is an
        isotropic ring. With $\sigma\approx150$, the whole critical shell grows
        in every orientation at once and becomes nonlinear in $t\approx0.04$,
        before mode selection can act. The clusters condense wherever the
        random superposition peaks. Rearranging them afterwards needs
        activated hops over large barriers (clusters about 190 times the mean
        density), so the disorder is frozen in.
  \item \textbf{Shallow quench: bcc crystal with twins and mosaic.}
        \begin{itemize}
          \item \emph{Crystal:} Bragg spots and $Q_6$ ten times the random
                floor.
          \item \emph{bcc:} the spot positions and angles, and a global $Q_4$
                near zero. A sizeable fcc or hcp domain keeps $Q_4\approx0.1$
                to 0.2.
          \item \emph{Not a single crystal:} $Q_6/q_6=0.5$, against 0.9 for one
                noisy crystal and 0.46 for two random grains. The twin domain
                and the sub-grains share the intensity, and the 25 clusters
                beyond the ideal 432 are defects.
        \end{itemize}
        With $\sigma\approx3.5$ only modes on the ring grow, and slowly, so the
        triads select the bcc $\langle110\rangle$ set while the amplitude is
        still small. The commensurate box-aligned set wins, and the clusters
        then condense onto a lattice that is already chosen. This is the field
        version of the rule that a deep quench gives a glass and a quench just
        past the transition nucleates a crystal.
\end{itemize}
The deep-quench data stop at $t=0.04$, against $t=0.17$ for the shallow quench,
so annealing time is not matched. But the deep quench showed no growth of
global order, while the shallow quench shows exactly the triad-selected bcc
set. So the onset mechanism, not elapsed time, is the likely cause.

\section{How to interpret the diagnostics}
\begin{center}
\begin{tabularx}{\textwidth}{@{}lL@{}}
\toprule
observation & meaning \\
\midrule
$S$ on the shell grows at rate $\approx\sigma$ & linear regime, far from $\Dc$ \\
growth rate $\gg\sigma$; onset nearly independent of $\sigma$ & triad (quadratic) growth near $\Dc$ \\
$\Lambda<1$, $S\approx1/(1-\Lambda)$ for a long time, then a sudden rise & nucleation of the crystal in the bistable window $\Dc<D<\Dsn$ \\
lattice survives a restart at $D>\Dc$ & first-order transition; melts only above $\Dsn$ \\
uniform shell in $S(\mathbf k)$, PR $\approx0.3$--$0.6$, $Q_6$ near $0.4/\sqrt N$ & amorphous or liquid packing of clusters \\
discrete spots, PR $\ll0.3$, $Q_6\gg0.4/\sqrt N$ & crystal \\
$Q_6/q_6\approx0.9$ & single crystal; lower values mean several grains \\
global $Q_4\approx0$ with spots & bcc (or a polycrystal); fcc and hcp keep $Q_4\approx0.1$--$0.2$ \\
spots at $60^\circ/90^\circ$, shells $1:\sqrt2:\sqrt3:2$ & bcc \\
8 spots at $70.5^\circ$, second shell at $1.155$ & fcc \\
spot sets related by $180^\circ$ about $\langle111\rangle$ & bcc $\Sigma3$ twins \\
\bottomrule
\end{tabularx}
\end{center}
For 3D bond order of bcc, use \code{--nn-cut} between its second and third
shells (0.2 here) or \code{--nn 14}. Twelve neighbours cut through its 8 + 6
first two shells.

\section{Caveats}
\begin{itemize}
  \item The amplitude equations are formally valid only for small amplitude.
        The cluster states are strongly nonlinear ($\phi_{\max}/\mu_0\approx6$
        to 190), so their agreement with the continuation (within 2--4\%) is
        partly fortuitous.
  \item The continuation compares one lattice, at one spacing (or a few), with
        the uniform state. It ignores coexistence, vacancies and multiple
        occupancy.
  \item The nucleation estimates treat noise as Gaussian fluctuations of
        independent modes. They predict the trend in
        $A_u/A_{\mathrm{noise}}$, not rates.
  \item In 3D at $128^3$ the clusters are only one or two cells wide, so
        resolution limits the cluster shapes and the low-$D$ free energies.
  \item The 3D spot analysis uses ten frames at $t\approx0.171$. Whether the
        twins and sub-grains anneal needs a longer series of $Q_6$ and of the
        brightest-12 share.
\end{itemize}

\section{Reproducing the results}\label{sec:repro}
From \code{exec/dean\_kow/interaction} (details in
\code{README\_crystal\_analysis.md}):
\begin{verbatim}
P=../../../python
# growth against linear theory, and nucleation above D_crit
python3 $P/sk_growth.py _new_2d_10_vvclose/plt_fv_* --num-part 327680 \
    --D 0.545 --sigma 5.61
python3 $P/sk_growth.py _new_2d_10_labove/plt_fv_* --num-part 163840 \
    --shell 8 --every 5
# amplitude-equation coefficients
python3 $P/amplitude_eqs.py --dim 2 --inputs inputs_fv --k 8 --dcrit 0.5472
python3 $P/amplitude_eqs.py --dim 3 --inputs inputs_fv_3d
# deterministic continuation
python3 $P/lattice_continuation.py --lattice hex --inputs inputs_fv --k 8 \
    --D 0.55 0.60 0.64 0.6475 0.65
python3 $P/lattice_continuation.py --lattice bcc --inputs inputs_fv_3d \
    --D 0.30 0.40 0.43 0.44 0.45 --tmax 0.06
# Bragg spots, bond order and reference values
python3 $P/sf_spots.py near_3d/plt_fv3d_* --every 10
python3 $P/cluster_diagnostics_3d.py near_3d/plt_fv3d_110900 --num-part 2e7 \
    --nn-cut 0.2 -o diag
python3 $P/order_reference.py --lattice bcc --cells 6 --nn-cut 0.2
\end{verbatim}

\begin{thebibliography}{9}
\bibitem{likos} C.~N.~Likos, A.~Lang, M.~Watzlawek and H.~L\"owen,
  \emph{Phys. Rev. E} \textbf{63}, 031206 (2001).
\bibitem{mladek} B.~M.~Mladek, D.~Gottwald, G.~Kahl, M.~Neumann and
  C.~N.~Likos, \emph{Phys. Rev. Lett.} \textbf{96}, 045701 (2006).
\bibitem{ciliberto} S.~Ciliberto, P.~Coullet, J.~Lega, E.~Pampaloni and
  C.~P\'erez-Garc\'ia, \emph{Phys. Rev. Lett.} \textbf{65}, 2370 (1990).
\bibitem{crossho} M.~C.~Cross and P.~C.~Hohenberg, \emph{Rev. Mod. Phys.}
  \textbf{65}, 851 (1993).
\bibitem{alexander} S.~Alexander and J.~McTague, \emph{Phys. Rev. Lett.}
  \textbf{41}, 702 (1978).
\end{thebibliography}

\end{document}
